Beveztés

\subsection{Biodiverzitási kockázatok számszerűsítése}

\initial{A} definícióból adódóan rögtön látszódik, hogy a biodiverzitás egy sokrétű fogalom, aminek a közgazdaságtani hatásait sok eltérő oldalról lehet megragadni. Az ismert ESG minősítők (MSCI, Bloomberg, LSEG Refinitiv) közül sokan publikálnak az E oszlopon belül környezeti (Nature) értékeket is, amiket mi is használni fogunk. A minősítőkön túl vannak vállalatok, akik speciálisabban foglalkoznak biodiverzitási kérdésekkel, nekik saját adatbázisaik is vannak.

Az LSEG Refinitiv platformról...

\subsection{Adatbázis}

A vizsgált befektetési univerzumunk az STOXX Europe 600 részvényindexre támoszkodik, amely az Eurozónán felül más európai országok részvényeit is tartalmazza. Az index több különböző devizában kereskedett eszközöket is tartalmaz, legnagyobb arányban az Euróban jegyzett részvények dominálnak, viszont emellett nagy részarányt képviselnek az angol fontban, valamint svájci frankban jegyzettek is.

A kutatásunkhoz használt biodiverzitási mutatók a következők:

\begin{table}[h!]
\centering
    \begin{tabular}{ | l | l |}
        \hline
        \textbf{Biodiverzitási mutató} & \textbf{Mutató jelentése} \\
        \hline
        Biodiversity Theme Score & aasgag \\
        Biodiversity Theme Materiality Level & aaaa \\
        Biodiversity Targets & aaa \\
        Biodiversity Impact Reduction & aaa \\
        Biodiversity Impact Reduction Score & aaa \\
        Biodiversity Risk Assessment & aaa \\
        Biodiversity Risk Management Integration & aaa \\
        Biodiversity Due Diligence & aaa \\
        \hline
    \end{tabular}
    \caption{Biodiverzitási mutatók}
\end{table}

Milyen időtávra vannak adatok...


Miket szűrtünk ki...

\subsection{Elsőkörös eredmények}

Iparágtól mennyire függ a mutató?
Időben hogyan változik? 2023-2025
